\section{Discussion}
\label{sec:discussion}

\paragraph{Practical guidance for method selection.}
Our results suggest that no single method universally dominates; the optimal choice depends on the
operational constraint.
When the primary constraint is \emph{compute}, MMR offers the most favorable accuracy--time
trade-off: its only overhead relative to Vanilla LoRA is the cost of sampling from the replay
buffer and computing loss on a slightly larger concatenated batch, which is negligible compared to
the forward-backward pass on current-task data.
When the primary constraint is \emph{data retention} (e.g., prior task data may not be stored due
to privacy regulations) while GPU memory is ample, O-LoRA is preferable because it enforces task
separation in parameter space without retaining any training examples; its higher wall-clock cost
is the trade-off.
When \emph{neither replay nor additional memory} for storing LoRA snapshots is feasible and the
practitioner has access to a held-out validation set for Fisher estimation, LoRA + EWC provides a
middle ground, though our results suggest its benefits are modest at the four-task scale studied
here.
Full fine-tuning should only be considered in single-task or offline multi-task settings; it
consistently incurs the most severe forgetting in our experiments despite achieving competitive
performance on the final task in the sequence.

\paragraph{Limitations.}
This study has several limitations that should temper the generality of its conclusions.
First, we evaluate on only four tasks drawn from a single benchmark suite (GLUE); tasks within
GLUE are relatively homogeneous in format and domain compared to real-world continual learning
deployments spanning scientific text, code generation, or dialogue.
Second, all tasks are classification problems cast as text generation via prompting; generative
tasks such as abstractive summarization or open-ended question answering may exhibit qualitatively
different forgetting dynamics, as the output space is unbounded and task interference may manifest
differently in the language modeling objective.
Third, both models evaluated are small ($\leq$4B parameters); at larger scales, richer
representations may be more amenable to orthogonalization (benefiting O-LoRA) or may exhibit
different forgetting rates due to increased model capacity.
Fourth, wall-clock measurements were taken on Colab's shared T4 infrastructure, which exhibits
variable GPU utilization; timing comparisons should therefore be treated as indicative rather than
definitive.

\paragraph{Threats to validity.}
Two additional threats deserve explicit mention.
\emph{Prompt format sensitivity} is a known confound for instruction-tuned models: slight
rephrasing of the task instruction can shift absolute accuracy by several points without changing
task content~\cite{luo2023empirical}.
We mitigate this by fixing a single prompt template per task (Section~\ref{sec:setup}), but
different template designs could alter the relative ranking of methods, particularly for tasks with
ambiguous class labels.
\emph{Task ordering effects} represent a second threat: prior work has shown that the order in
which tasks are presented can substantially change forgetting patterns, as earlier tasks in the
sequence have more gradient updates applied on top of them.
We evaluate a single fixed sequence (SST-2\,$\to$\,MRPC\,$\to$\,CoLA\,$\to$\,MNLI) chosen to
maximize linguistic diversity, but a full permutation analysis across all $K! = 24$ orderings is
beyond the compute budget of this study and is left as future work.
